\documentclass[../main]{subfiles}
\begin{document}


\chapter*{Introduction}
%This is the background and larger scope of the paper: 
\section*{Methane and Climate Sensing}
Climate Sensing 
    




Need for model
    mixing rations and Lab spectroscopy
        examples of uses of lab spectroscopy in history and importance in climate policy/action




In addressing climate change, we must frequently take measurements of greenhouse gasses in Earth’s atmosphere. 
To measure the amount of methane in the atmosphere, we must first have a model of the intensity of the absorption spectrum of 
methane at different temperatures and pressures. Uncertainties in these models, particularly in lineshape functions, 
propagate to uncertainties in greenhouse gas measurements. 
\section*{Problems in Modeling Methane's Behavior}
Spectroscopic models of greenhouse gasses used in measurement of the atmosphere typically neglect 
accounting for collisional effects that lead to broadening and narrowing effects of the line shape of the spectra 
produced by greenhouse gasses. 
\section*{Ro-Vibrational Transitions and HITRAN Data}




\section*{Contribution Made by this Work}
We take our measurements with a novel digitally controlled IR laser spectrometer created using the FP4209 InstaTune laser module 
with the intent to characterize the effects of pressure on the line shape function used in modeling the predicted absorption spectrum of methane. 
The FP4209 scans over a IR wavelength range of 1627.5000 – 1672.4955 nm, which corresponds very well with the ro-vibrational absorption spectrum of methane. 
Interpolating DAC values and scanning through the range of wavelengths produced by this laser allows us to construct a quasi-continous profile of 
the spectrum of the intensity of light absorbed by the methane gas. 

This work allows us to take new measurements of the properties of methane gas in a controlled environment, 
to better understand the spectra we observe in the atmosphere, and to more precisely describe the behavior 
and amount of methane in our atmosphere. 

\end{document}
